\section{Where the premium lives}
\label{sec:deamwhere}

We know the premium biases the naive reading.  Before removing it, it
pays to know where it is --- because over much of the strike--maturity
plane it is \emph{exactly zero}, and the theory that says so is short
enough to carry in full.

First the formal object.  A \emph{stopping time} is a rule for
choosing when to act that may use only what has happened so far ---
``exercise the first time the stock drops below $80$'' is a stopping
time; ``exercise at the low of the path'' is not, because recognizing
the low requires seeing the future.  Write $\mathcal{T}$ for the set
of all stopping times $\vartheta$ taking values in $[0,t]$.  The
American value is the value of the \emph{best plan}:
\begin{equation}
A=\sup_{\vartheta\in\mathcal{T}}
\E\!\left[e^{-r\vartheta}\,
\text{payoff}(S_\vartheta)\right],
\label{eq:deamstop}
\end{equation}
with the payoff $(S_\vartheta-K)^+$ for a call and $(K-S_\vartheta)^+$
for a put.  The expectation is under the risk-neutral law ---
Chapter~2's pricing measure, and the only fact about it this section
uses is that it grows the spot at the carry, the rate the forward
implies.  Discounting runs at the rate $r$ on the calendar clock $t$,
as all financing has since Chapter~6 \citep{PeskirShiryaev2006}.

\begin{proposition}[Dominance, both ways]
\label{prop:deamdom}
$A\ge E$, and $A\ge$ intrinsic; hence the premium $A-E$ is never
negative, and an American quote never sits below its exercise floor.
\end{proposition}

\begin{proof}
The plans ``always wait to expiry'' ($\vartheta\equiv t$) and
``exercise immediately'' ($\vartheta\equiv0$) are both stopping times.
The first values the European twin; the second values the intrinsic.
A supremum is at least every member of its feasible set.
\end{proof}

One line, but it already fixes two working conventions.  Because the
true premium cannot be negative, a negative \emph{estimate} of it can
be clamped to zero with a clear conscience (\cref{sec:deamsub}).  And
an American quote \emph{below} intrinsic is not a puzzle about
premiums --- it is a broken price, to be discarded, not interpreted.
The substantive question is when the inequality $A\ge E$ is an
equality: when is the extra right worth nothing?

\begin{theorem}[Merton: no dividends, no call premium]
\label{thm:deammerton}
For a call on an underlying that pays no dividends, with $r\ge0$,
early exercise is never strictly better than waiting: $A=E$
\citep{Merton1973}.
\end{theorem}

\begin{proof}
Consider the option at any interior moment, with spot $s$ and time $u>0$
still to run, and compare exercising now against one specific
alternative plan: waiting to the end.  Waiting is worth the European
call value, and three displayed steps bound it from below.  The hinge
function $x\mapsto(x-K)^+$ is convex, so by Jensen's inequality ---
the average of a convex function is at least the function of the
average ---
\[
e^{-ru}\,\E\big[(S_u-K)^+ \,\big|\, S_0=s\big]
\;\ge\;
e^{-ru}\big(\E[S_u\,|\,S_0=s]-K\big)^{+}.
\]
With no dividends, the risk-neutral spot grows at the full rate $r$
(nothing leaks out of the stock): $\E[S_u\,|\,S_0=s]=s\,e^{ru}$.
Substituting,
\[
e^{-ru}\big(s\,e^{ru}-K\big)^{+}
=\big(s-Ke^{-ru}\big)^{+}
\;\ge\; s-K,
\]
where the last inequality holds because $Ke^{-ru}\le K$, and it is
\emph{strict} whenever $r>0$ and $u>0$.  So at every interior moment,
the value of waiting strictly exceeds $s-K$, which is everything
exercising now can pay.  Exercising early forfeits value at every
opportunity; the best plan in \eqref{eq:deamstop} is
$\vartheta\equiv t$, and that plan is the European contract.
\end{proof}

In words a colleague could repeat: exercising a call early hands over
the strike $K$ today when the contract only ever requires handing it
over at expiry --- giving up the interest on the strike --- and
collects a stock that, paying no dividends, was growing at the full
rate anyway.  There is nothing to gain and interest to lose.  The
right is legally present and economically worthless.

Puts are different, and the same interest-on-the-strike logic now
works \emph{for} exercising.

\begin{proposition}[Deep puts must carry a premium]
\label{prop:deamput}
For a put with $r>0$, the premium is strictly positive deep enough in
the money: whenever $S<K\,(1-e^{-rt})$, exercising now strictly beats
waiting to expiry, so $A>E$.
\end{proposition}

\begin{proof}
A put can never pay more than $K$, and it pays at expiry, so its
European value is capped by the discounted strike: $E\le Ke^{-rt}$.
Deep enough in the money, the intrinsic value exceeds that cap:
\[
K-S\;>\;Ke^{-rt}
\quad\Longleftrightarrow\quad
S\;<\;K\big(1-e^{-rt}\big).
\]
There, $A\ge K-S>Ke^{-rt}\ge E$: the exercise-now plan strictly beats
the European plan, and the premium $A-E$ is positive.
\end{proof}

Read the trade-off plainly.  The put holder deep in the money is
essentially sure to receive $K$; receiving it \emph{now} rather than
at expiry earns the interest.  Waiting to the end pays only the
option's remaining \emph{time value} --- its price in excess of
intrinsic --- which deep in the money shrinks toward zero while the
interest does not.  The threshold in the proposition is
crude --- it only proves the exercise region is nonempty --- but the
mechanism is exact, and it has a mirror image: a \emph{dividend} is
value the \emph{call} holder collects only by exercising before the
stock goes ex-dividend, so dividends (or any carry $q_d$ exceeding
$r$) open an exercise region on deep in-the-money calls.
\Cref{sec:deamdivs} prices that mechanism properly.

\Cref{fig:deammap} draws the whole geography, computed by the tree of
\cref{sec:deamtree} at a fixed $25\%$ volatility.  Panel~(a) maps the
put premium $A-E$ in dollars over strike (as a percentage of spot) and
maturity, at $r=4\%$ with no dividends.  Two features to notice.
The premium concentrates deep in the money and grows with maturity ---
more calendar for the interest to accrue over --- reaching
\$\MacDeamMapPutMaxDollars\ per hundred dollars of spot at the top
right.  And the orange edge marks something stronger than a large
premium: to its right the price \emph{equals} intrinsic, to the
plateau tolerance stated in appendix~\ref{app:deamprotocol} --- at
half a year, every strike beyond $\MacDeamMapPutBdryPct\%$ of spot.
In that dead zone --- the intrinsic \emph{plateau}, and the chapter
will use both names: the price sits flat on its floor there --- the
quote is the exercise floor and nothing else: no volatility enters
it, so no volatility can be read out of it.
\Cref{sec:deamsub} shows how the inversion recognizes the zone and
refuses it honestly.  Panel~(b) runs the mirror image --- calls under
the running example's $6\%$ dividend yield against a $2\%$ rate ---
and the picture flips left to right: the premium sits on deep
in-the-money calls, and the dead zone opens below
$\MacDeamMapCallBdryPct\%$ of spot.  One number is deliberately absent
from the figure: on the same grid with the dividend yield set to zero,
the largest call premium anywhere is
$\MacDeamMapMertonMaxDollars$ --- not small, \emph{zero}: at every
node of every tree the waiting value exceeded intrinsic, the exercise
comparison never once engaged, and the American and European
recursions returned identical arrays.  That is \cref{thm:deammerton},
measured.

\begin{figure}[htbp]
\centering
\paperfig{\textwidth}{fig_deam_map.pdf}
\caption{The premium $A-E$ (dollars per 100 spot) over the quote
plane, at a fixed $25\%$ volatility.  (a)~Puts at $r=4\%$, no
dividends: the premium loads on deep in-the-money strikes and long
maturities; right of the orange edge the price equals intrinsic ---
at $t=0.5$, beyond $\MacDeamMapPutBdryPct\%$ of spot --- and carries
no volatility information.  (b)~Calls under a $6\%$ dividend yield
against a $2\%$ rate: the mirror image, dead zone below
$\MacDeamMapCallBdryPct\%$ of spot.  With the dividend yield set to
zero the call premium is exactly zero everywhere on the grid:
Merton's theorem, measured.}
\label{fig:deammap}
\end{figure}

We now know the premium's address: deep in-the-money puts under
positive rates, deep in-the-money calls around dividends, growing
with maturity, zero for no-dividend calls, and --- past the plateau
--- replaced by a dead zone where the quote has nothing to say about
volatility.  What we lack is a machine that prices the best-plan value
\eqref{eq:deamstop} at will.  Building one small enough to run by
hand is the next section's job.
